% page 630 of the facsimile (image p-629 of the scan)
% block 154.9 x 254.9 mm ; left margin 8.0 mm
\pgc{-12.700mm}{0.000mm}{
\l{}
\l{}
\l{\sou{votar}. (\sou{netrans., pri, por, kontre}). (\sou{Aludante asemblitaro}}
\l{\cel{3}en\cel{1}\sou{qua la rezolvi esas agenda segun la opiniono da}\cel{1}la}
\l{\cel{3}\sou{majorit}ato - ed\cel{1}\sou{aludante singla membro}). Depozar billeto,}
\l{\cel{3}bulo, e c. qua indikas ke (ta membro) adoptas o ne-adop-}
\l{\cel{3}tas ta o ca propozajo. - DEFIRS.}
\l{}
\l{\sou{vovar}.\cel{2}(trans.) Promisar a la deajo exekutor ta o ca tasko,}
\l{\cel{3}verko meritoza, se lu exaucos la demando quan onu agas a}
\l{\cel{3}lu. - EF.}
\l{}
\l{\sou{voyo}. I. La spaco parkurenda por irar de loko ad altra loko,}
\l{\cel{3}equipita por la cirkulado di la veturi e di la pediranti.}
\l{\cel{3}- II. La spaco parkurita da la vildo quan onu persequas}
\l{\cel{3}e quan revelas la traci quin lu lasis. - F.}
\l{}
\l{\sou{voyajar}. (\sou{netrans., de, ad)}. Diplasar su, parkurante voyo}
\l{\cel{3}plu o min longa, por irar til altra urbo, lando o konti-}
\l{\cel{3}nento (o mem planeto). - EFIS.}
\l{}
\l{vu.\cel{2}(\sou{gram.}) Pronomo personala qua, dum nefamiliareso ye ulu,}
\l{\cel{3}remplasas \textquotedbl{}tu\textquotedbl{}. - FIR.}
\l{}
\l{\sou{vulgara}.\cel{2}I. Qua admisesas, uzesas da la homi ordinara. - II.}
\l{\cel{3}Quan nulo distingas de la homi ordinara.- DEFIRS.}
\l{}
\l{\sou{vulgato}. (\sou{religio kristana}). Tradukuro Latina de la texto}
\l{\cel{3}Hebrea di biblo, agnoskata kom konforma a la kanoni di}
\l{\cel{3}le eklezio katolika. - DEFIRS.}
\l{}
\l{\sou{vulkanala}. (\sou{geol.)}\cel{2}Qua imputas a fairo la formaco di la}
\l{\cel{3}Terglobo. -}
\l{}
\l{\sou{vulkanizar}. (trans.).(\sou{tekn.})\cel{2}Preparar (kauchuko) per imer-}
\l{\cel{3}sar lu en sulfo-fuzita por igar lu ne apta influesar da}
\l{\cel{3}la fluktui di temperaturo e di vetero. - DEFIRS.}
\l{}
\l{\sou{vulnerario}. (\sou{bot.)}\cel{3}Planto leguminosa, kun flori flava, re-}
\l{\cel{3}putata kom apta remediar la vunduri, la plagi. - L.\cel{1}\sou{vul}-}
\l{\cel{3}\sou{neraria rustica}. - FIS.}
\l{}
\l{\sou{vulto}.\cel{2}Konstrukturo di qua la suprajo esas kurvatra, ek}
\l{\cel{3}asembluro de quadri infre-kono-trunka qui inter-apogas.}
\l{\cel{3}- EFI.}
\l{}
\l{\sou{vulturo}. I. (\sou{zool.}) Granda raptucelo, kun kapo e kolo nuda,}
\l{\cel{3}garnisita ye nur lanugo, sen plumi. - L.\cel{1}\sou{vultur}.-\cel{2}II.}
\l{\cel{3}(\sou{metaf}.) La persono da qua la febli esas viktima. - EFI.}
\l{}
\l{\sou{vulvo}. (\sou{anat.})\cel{2}Orifico di la organi genitala che la femino.}
\l{\cel{3}- DEFIRS.}
\l{}
\l{\sou{vundar}. (\sou{trans.}) (ye).-\cel{2}I. Frapar per stroko qua lezas, de-}
\l{\cel{3}trimentas la organismo. - II. (\sou{aludante objekti}) Lezar}
\l{\cel{3}parto di korpo per frotado, aracho, e c. - III. Afektar}
\l{\cel{3}penigive (organo di la sensi). -\cel{2}DE.}
\l{}
\l{\cel{30}----}
}
